\section{References}
\label{sec:references}%

The following references were consulted and used during the development of this Requirements Analysis and Specification Document:

\begin{enumerate}
    \item \textbf{IEEE Standard} - IEEE Recommended Practice for Software Requirements Specifications.
    \item \textbf{Specification for the 2025/2026 Requirement Engineering and Design Project} - Project assignment document.
    \item \textbf{Software Engineering 2 Course Materials}, Academic Year 2024-2025 - Course slides available on the WeBeep page.
    \item \textbf{Previous year student project} - Available on the WeBeep page for reference and guidance.
    \item \textbf{General Data Protection Regulation (GDPR)} - Regulation on data protection and privacy.
\end{enumerate}

\section{Used Tools}
\label{sec:used_tools}%

The following tools and technologies were utilized during the development of this document:

\begin{enumerate}
    \item \textbf{GitHub} - Version control and collaborative development platform for managing the project repository.
    \item \textbf{\LaTeX{}} - Document preparation system used for writing and formatting this document, with Overleaf and Neovim as editing environments.
    \item \textbf{Visual Studio Code} - Code editor used for writing and editing LaTeX files and other project documentation.
    \item \textbf{Draw.io} - Diagramming tool for creating UML diagrams and flowcharts.
    \item \textbf{Alloy Analyzer} - Formal specification language and analyzer for model checking and validation.
    \item \textbf{Google Documents} - Collaborative platform for brainstorming, shared notes, and draft document preparation.
\end{enumerate}

\subsection{Use of Generative AI Tools}
\label{subsec:ai_tools}%

Generative AI tools were used as support tools for specific, well-defined tasks during the development of this document. The following section describes the tools used, the inputs provided, and the outputs obtained.

\subsubsection*{Tool Used}
\textbf{ChatGPT} - Large Language Model chatbot used for clarification of concepts, text rewriting, and assistance in structuring content.

\subsubsection*{Inputs and Outputs}

The following queries were used during the document preparation:

\begin{enumerate}
    \item \textbf{Query:} "Hello, I'm addressing the scope of a project for the RASD document (requirement analysis and specification document) and I can't differentiate the world phenomena and the shared phenomena in a project. Give me some examples and I'll try to differentiate by myself first"
    
    \textbf{Output:} The tool provided examples and explanations that helped distinguish between phenomena that exist independently in the world versus those that are shared between the system and its environment.
    
    \item \textbf{Query:} "Help me rewrite more clearly these user scenarios: [paste previous user scenarios]"
    
    \textbf{Output:} The tool provided rewritten versions of the scenarios with improved clarity, better structure, and more precise language.
    
    \item \textbf{Query:} "How can I rewrite more clearly the domain assumptions and regulatory policies? [paste previous domain assumptions and regulatory policies]"
    
    \textbf{Output:} The tool suggested improvements in structure, formatting, and language clarity for the domain assumptions and regulatory policies.
    
    \item \textbf{Query:} "Given the two screenshots attached, generate some mockup UI designs that could help inspire the rest of the frontend design for our Best Bike Paths mobile application. Consider the style, layout, and design patterns shown in the screenshots and create additional mockup suggestions that follow a similar aesthetic and user experience approach."
    
    \textbf{Output:} The tool generated mockup UI designs that provided inspiration for the frontend design, suggesting layout patterns, color schemes, and interface elements consistent with the style shown in the provided screenshots.
\end{enumerate}

\vspace{0.3cm}
\textit{Note: Additional queries were used during the document preparation, but as they were redundant with the queries presented above, they have not been included in this document.}

\subsubsection*{Verification and Refinement Process}

All outputs from the generative AI tool were subject to the following verification process:
\begin{itemize}
    \item \textbf{Technical Accuracy:} All suggestions were cross-referenced with course materials, IEEE standards, and the project assignment specifications.
    \item \textbf{Content Validation:} Generated content was reviewed to ensure it accurately reflected the Best Bike Paths system requirements and did not introduce incorrect or misleading information.
    \item \textbf{Integration Review:} Before integration, all AI-generated suggestions were evaluated by both team members to ensure consistency with the document's overall structure and style.
    \item \textbf{Final Verification:} All integrated content was reviewed in the context of the complete document to ensure coherence and accuracy.
\end{itemize}

\subsubsection*{Responsibility Statement}

The authors take full responsibility for all content in this document. The generative AI tool was used solely as a support tool for clarification, text improvement, and structural guidance. All final content, including text, diagrams, and specifications, was created, verified, and approved by the authors, who are fully capable of explaining and defending every aspect of the work presented in this document.
